\section{Non-Rapid Eye Movement (NREM) Sleep}\label{sec:nrem}

NREM sleep constitutes the majority of sleep time (about 75--80~\%) and is the state we first enter. It is itself a progression through three stages of increasing depth:

\begin{itemize}
    \item \textbf{N1 (Light Sleep):} This is the brief, transitional stage between wakefulness and sleep. Muscle activity slows, and it's common to experience hypnagogic jerks or the sensation of falling. This stage is also associated with hypnagogic hallucinations, which are explored in Chapter~\ref{ch:phenomenology}.
    
    \item \textbf{N2 (Stable Sleep):} This is the first ``true'' stage of sleep, where you spend the largest percentage of your night. Brain waves slow down, punctuated by sudden bursts of activity known as \textit{sleep spindles} and \textit{K-complexes}. These are thought to be mechanisms for memory consolidation and suppressing external stimuli~\cite{fuller2006}.
    
    \item \textbf{N3 (Deep Sleep / Slow-Wave Sleep):} This is the deepest, most restorative stage of sleep, characterized by very slow, high-amplitude \textit{delta waves}. It is most prevalent in the first third of the night. During N3, the body performs critical maintenance: the glymphatic system clears metabolic waste from the brain (see Chapter~\ref{ch:functions}), and the pituitary gland releases human growth hormone for tissue repair~\cite{benington1999}.
\end{itemize}
